% !TeX program = xelatex
\documentclass[11pt,oneside]{article}
\usepackage{ProblemsetStyle}

\hypersetup{
    pdftitle={20260913 Problem set - Solutions},
    pdfauthor={김정우},
    pdfsubject={Science Quiz mathematics practice problem set solutions},
    pdfcreator={XeLaTeX}
}

\begin{document}

\problemsettitle{20260913 Problem set}

\begin{problemsolution}{1}
Anthropic이 2026년 9월 4일 공개한 formalization의 대상은 Fermat's Last Theorem이다. Claude는 Lean으로 기존 증명을 end-to-end formalize하여 computer-checked proof를 완성했다.

따라서 정답은
\[
    \boxed{\text{Fermat's Last Theorem}}
\]
이다.
\end{problemsolution}

\begin{problemsolution}{2}
$101$은 홀수 소수이므로 $37\not\equiv0\pmod{101}$에 대해 해의 개수는 $37$이 quadratic residue인지에 따라 $0$ 또는 $2$이다. Quadratic reciprocity를 쓰면
\[
    \left(\frac{37}{101}\right)
    =
    \left(\frac{101}{37}\right)
    =
    \left(\frac{27}{37}\right)
    =
    \left(\frac{3}{37}\right).
\]
여기서 $37\equiv1\pmod{12}$이므로
\[
    \left(\frac{3}{37}\right)=1.
\]
따라서 $37$은 modulo $101$에서 quadratic residue이고, 서로 다른 두 square root를 갖는다. 실제로
\[
    21^2\equiv80^2\equiv37\pmod{101}.
\]
그러므로 residue class의 개수는
\[
    \boxed{2}.
\]
\end{problemsolution}

\begin{problemsolution}{3}
Stationary distribution에서 상태 $A$일 확률을 $p$라 하자. 한 단계 뒤에도 이 확률이 같아야 하므로
\[
    p
    =
    \frac23p+\frac14(1-p).
\]
따라서
\[
    \frac{7}{12}p=\frac14,
    \qquad
    p=\frac37.
\]
그러므로 정답은
\[
    \boxed{\frac37}.
\]
\end{problemsolution}

\begin{problemsolution}{4}
University of Oxford가 발표한 Andrew Wiles의 후임 Regius Professor of Mathematics는 James Maynard이다. Maynard는 analytic number theory, 특히 prime gaps와 sieve methods에 관한 업적으로 2022 Fields Medal을 받았다.

따라서 정답은
\[
    \boxed{\text{James Maynard}}.
\]
\end{problemsolution}

\begin{problemsolution}{5}
$P(x,y)=-y^2$, $Q(x,y)=x^2$라 두면 Green's theorem에 의해
\[
    \oint_C P\,\dd x+Q\,\dd y
    =
    \iint_D
    \left(
        \frac{\partial Q}{\partial x}
        -
        \frac{\partial P}{\partial y}
    \right)
    \dd A,
\]
여기서
\[
    D=\{(x,y)\in\R^2\mid x\geq0,\ y\geq0,\ x+y\leq2\}.
\]
따라서
\[
    \begin{aligned}
    \oint_C -y^2\,\dd x+x^2\,\dd y
    &=
    \int_0^2\int_0^{2-x}2(x+y)\,\dd y\,\dd x\\
    &=
    \int_0^2(4-x^2)\,\dd x
    =\frac{16}{3}.
    \end{aligned}
\]
그러므로
\[
    \boxed{\frac{16}{3}}.
\]
\end{problemsolution}

\begin{problemsolution}{6}
첫 번째 column은 $\F_3^2$의 nonzero vector이면 되므로 선택지는
\[
    3^2-1=8
\]
개이다. 첫 번째 column을 고정하면 그 span에는 $3$개의 vector가 있으므로, 두 번째 column은 그 span 밖의
\[
    3^2-3=6
\]
개 중 하나여야 한다. 따라서
\[
    |\GL_2(\F_3)|
    =(3^2-1)(3^2-3)
    =8\cdot6
    =\boxed{48}.
\]
\end{problemsolution}

\begin{problemsolution}{7}
ICTP와 IMU가 발표한 2025 Ramanujan Prize 수상자는 Claudio Muñoz이다. 수상 사유는 dispersive partial differential equations의 long-time behavior, 특히 soliton의 asymptotic stability와 multi-soliton dynamics에 대한 기여였다.

따라서 정답은
\[
    \boxed{\text{Claudio Mu\~noz}}.
\]
\end{problemsolution}

\newpage

\begin{problemsolution}{8}
Shoelace formula로 quadrilateral의 area는
\[
    A
    =
    \frac12(5\cdot4+3\cdot3)
    =
    \frac{29}{2}.
\]
각 변 위의 lattice segment 수를 더하면 boundary lattice point의 개수는
\[
    B
    =
    \gcd(5,0)+\gcd(2,4)+\gcd(3,1)+\gcd(0,3)
    =5+2+1+3
    =11.
\]
Pick's theorem
\[
    A=I+\frac{B}{2}-1
\]
을 적용하면
\[
    I
    =
    \frac{29}{2}-\frac{11}{2}+1
    =10.
\]
따라서 내부 lattice point의 개수는
\[
    \boxed{10}.
\]
\end{problemsolution}

\begin{problemsolution}{9}
Clay Mathematics Institute는 Yu Deng과 Zaher Hani가 hard-sphere system의 microscopic dynamics에서 출발하여 long times에 대해 Boltzmann equation을 rigorous하게 유도한 업적을 수상 사유로 들었다.

따라서 정답은
\[
    \boxed{\text{Boltzmann equation}}.
\]
\end{problemsolution}

\begin{problemsolution}{10}
$|z|=1$에서
\[
    |z^5+1|
    \leq
    |z|^5+1
    =2
    <5
    =|5z^2|.
\]
따라서 Rouch\'e's theorem에 의해
\[
    z^5+5z^2+1
\]
과 $5z^2$는 open unit disk 안에 같은 개수의 zero를 갖는다. $5z^2$는 $z=0$에서 multiplicity $2$의 zero를 가지므로, 주어진 다항식의 zero 개수도
\[
    \boxed{2}
\]
이다.
\end{problemsolution}

\end{document}
